\section{Literature Review}

% =============================================================================
% Overview
% =============================================================================
\begin{frame}{Literature Review: Overview}
    \begin{block}{Focus Areas}
        Critical review of literature concerning:
        \begin{itemize}
            \item Online assessment systems evolution (COVID-19 pandemic impact)
            \item Online proctoring and integrity assurance
            \item AI-driven Automated Essay Scoring (AES)
            \item Large Language Models (LLMs) in assessment
            \item Natural Language Processing (NLP) advancements
        \end{itemize}
    \end{block}
    
    \begin{block}{Key References}
        \begin{itemize}
            \item Topuz (2022): Systematic review of 22 online assessment systems
            \item Hussein (2020): Commercial proctoring tools evaluation
            \item Mizumoto \& Eguchi (2023): GPT-3 for AES
            \item Schneider (2024): GPT-4 for short answer grading
        \end{itemize}
    \end{block}
\end{frame}

% =============================================================================
% Evolution Phases
% =============================================================================
\begin{frame}{Evolution of Online Assessment Systems}
    \begin{table}[]
        \centering
        \small
        \begin{tabular}{l p{8cm}}
            \toprule
            \textbf{Phase} & \textbf{Key Characteristics} \\
            \midrule
            \textbf{Phase 1} & \textbf{Manual Online (Pre-2020):} Basic digital delivery, webcam monitoring, \\
            (Pre-2020) & browser lockdowns. Limited real-world deployment. \\
            \midrule
            \textbf{Phase 2} & \textbf{Emergency Remote Teaching (2020-2021):} Rapid adoption, \\
            (2020-2021) & desktop-focused, facial recognition, semi-automatic monitoring, \\
            & LMS integration (Moodle). Infrastructure gaps exposed. \\
            \midrule
            \textbf{Phase 3} & \textbf{AI Integration (2021-2024):} Machine learning for cheating \\
            (2021-2024) & detection, real-time monitoring, deep learning for behavior analysis. \\
            & Issues: false positives, bias concerns. \\
            \midrule
            \textbf{Phase 4} & \textbf{User-Centric \& Ethics (2024-2025):} Privacy focus, GDPR \\
            (2024-2025) & compliance, hybrid models, human oversight, equity considerations. \\
            \midrule
            \textbf{Phase 5} & \textbf{LLM-Driven Scoring (2023-2026):} GPT-3/4, ChatGPT, Llama \\
            (2023-2026) & for AES and ASAG. Rubric-based evaluation, feedback generation. \\
            \bottomrule
        \end{tabular}
    \end{table}
\end{frame}

% =============================================================================
% Study Related Works
% =============================================================================
\begin{frame}{Key Related Works (1/3)}
    \begin{block}{Topuz et al. (2022) - Systematic Review}
        \begin{itemize}
            \item Reviewed 22 online assessment systems during Emergency Remote Teaching
            \item \textbf{Major features identified:}
            \begin{itemize}
                \item Multi-platform compatibility
                \item Enhanced security (biometric authentication)
                \item Automated proctoring
            \end{itemize}
            \item Commercial systems: Examity, ProctorU, Safe Exam Browser
            \item Multimodal authentication: webcam, keystroke analysis, facial recognition
        \end{itemize}
    \end{block}
\end{frame}

\begin{frame}{Key Related Works (2/3)}
    \begin{block}{Kurniati et al. (2023) - Indonesian Context}
        \begin{itemize}
            \item Challenges: Infrastructure limitations, academic dishonesty, inconsistent discipline
            \item \textbf{Adaptive strategies:} Shift to subjective assessments, home visits, shared mobile access
        \end{itemize}
    \end{block}
    
    \begin{block}{Hussein (2020) \& Arno (2021)}
        \begin{itemize}
            \item \textbf{Hussein:} Cost, accessibility, ethical concerns as barriers
            \item \textbf{Arno:} Compared 29 commercial systems
            \item \textbf{Conclusion:} Hybrid approaches (automated + human) optimal for integrity \& UX
        \end{itemize}
    \end{block}
\end{frame}

\begin{frame}{Key Related Works (3/3)}
    \begin{block}{AI-Based Proctoring}
        \begin{itemize}
            \item \textbf{Jia \& He (2021):} IOPS - facial, eye, voice monitoring. Reduced cheating in university exams
            \item \textbf{Rahim (2020):} Guidelines for effective online assessment during ERT
            \item \textbf{Ali (2025):} Phased implementation in medical education. Mock exams improved confidence
        \end{itemize}
    \end{block}
    
    \begin{block}{LLM-Based Grading}
        \begin{itemize}
            \item \textbf{Mizumoto \& Eguchi (2023):} GPT-3 for TOEFL essays (QWK $\sim$0.7)
            \item \textbf{Schneider (2024):} GPT-4 for short answers - reliable on rubric-based benchmarks
            \item \textbf{Mansour (2024):} LLMs for proficiency scoring with trait-based feedback
            \item \textbf{Kundu (2024):} LLMs harsher on nuanced responses, useful as assistants
        \end{itemize}
    \end{block}
\end{frame}

% =============================================================================
% Milestones
% =============================================================================
\begin{frame}{Milestones: Phase 1 \& 2}
    \begin{block}{Phase 1: Manual Online Assessments (Pre-2020)}
        \textbf{Strengths:}
        \begin{itemize}
            \item Clear conceptual frameworks for design and security
            \item Authentication classification (knowledge, possession, biometric, contextual)
        \end{itemize}
        \textbf{Limitations:}
        \begin{itemize}
            \item Largely descriptive, limited real-world testing
            \item Minimal emphasis on scalability or low-bandwidth environments
        \end{itemize}
    \end{block}
    
    \begin{block}{Phase 2: Emergency Remote Teaching (2020-2021)}
        \textbf{Strengths:}
        \begin{itemize}
            \item Comprehensive overview of pandemic-driven solutions
            \item Desktop-focused, pre-exam facial recognition, LMS integration
        \end{itemize}
        \textbf{Limitations:}
        \begin{itemize}
            \item Self-reported data, overlooks infrastructure constraints in developing regions
        \end{itemize}
    \end{block}
\end{frame}

\begin{frame}{Milestones: Phase 3, 4 \& 5}
    \begin{block}{Phase 3: AI \& Automation (2021-2024)}
        \textbf{Strengths:} Proactive, automated integrity enforcement \\
        \textbf{Limitations:} Requires high-quality data, stable internet; bias, fairness, privacy concerns
    \end{block}
    
    \begin{block}{Phase 4: User-Centric Design (2024-2025)}
        \textbf{Strengths:} Ethics, privacy, GDPR compliance, human oversight \\
        \textbf{Limitations:} Early stages, limited empirical evidence, Western-centric
    \end{block}
    
    \begin{block}{Phase 5: LLM-Driven Scoring (2023-2026)}
        \textbf{Strengths:} Scalability, personalization, instant feedback (QWK 0.68-0.88) \\
        \textbf{Limitations:} Prompt design sensitivity, training biases, variable correlations, limited validation in resource-constrained environments
    \end{block}
\end{frame}

% =============================================================================
% Gaps Identified
% =============================================================================
\begin{frame}{Critical Gaps Identified (1/3)}
    \begin{alertblock}{1. Poor Accessibility for Students with Disabilities}
        \begin{itemize}
            \item Restrictive security measures (browser lockdowns, hidden taskbars)
            \item Exclude users dependent on assistive technologies
            \item Lack of Universal Design for Learning (UDL) principles
            \item Inconsistent accessibility policies across institutions
            \item \textbf{Design bias:} Security \& standardization over inclusivity
        \end{itemize}
    \end{alertblock}
\end{frame}

\begin{frame}{Critical Gaps Identified (2/3)}
    \begin{alertblock}{2. The ``Security vs. Connectivity'' Gap}
        \begin{itemize}
            \item Secure proctoring requires continuous high-speed internet
            \item Weak digital infrastructure in developing regions
            \item Indonesian teachers abandoned video-based testing (Kurniati 2023)
            \item 57\% of Pakistani medical students experienced internet instability (Ali 2025)
            \item Commercial tools require intensive data streaming, high computing power
            \item \textbf{Missing:} Adaptive, low-bandwidth monitoring mechanisms
        \end{itemize}
    \end{alertblock}
    
    \begin{alertblock}{3. The ``Intelligent Proctoring'' Gap}
        \begin{itemize}
            \item Polarization: Simple unproctored tools vs. costly AI-driven systems
            \item High-end platforms financially inaccessible to smaller universities
            \item \textbf{Missing:} Accessible ``Passive AI'' solutions (tab-switch detection, periodic image capture, local face verification)
        \end{itemize}
    \end{alertblock}
\end{frame}

\begin{frame}{Critical Gaps Identified (3/3)}
    \begin{alertblock}{4. The ``Local Customization'' Gap}
        \begin{itemize}
            \item Technologies developed for Western institutions
            \item Seldom consider local privacy laws, language interfaces, workflows
            \item Nearly all systems hosted abroad
            \item Difficulties adapting foreign platforms to local realities
            \item \textbf{Missing:} Indigenous, region-specific frameworks compliant with national data protection (e.g., Ethiopia's Personal Data Protection Proclamation)
        \end{itemize}
    \end{alertblock}
    
    \begin{alertblock}{5. ASAG \& AES Gaps}
        \begin{itemize}
            \item \textbf{ASAG:} Variable GPT-4 performance, inconsistent domain-specific scoring
            \item \textbf{AES:} Variable human-LLM agreement, harsher on nuanced essays, limited non-English/technical domain adaptation
        \end{itemize}
    \end{alertblock}
\end{frame}

% =============================================================================
% Lessons Learned
% =============================================================================
\begin{frame}{Lessons Learned from Literature}
    \begin{enumerate}
        \item \textbf{Infrastructure \& Inclusivity Matter}
        \begin{itemize}
            \item Digital divide's role in exam integrity
            \item Flexible system design for varying connectivity
        \end{itemize}
        
        \item \textbf{Ethical \& Data-Conscious Design}
        \begin{itemize}
            \item Transparent data handling, user-centric approaches
        \end{itemize}
        
        \item \textbf{Integrity Cannot Rely on Trust Alone}
        \begin{itemize}
            \item Technical methods essential (tab-switching, face detection)
            \item Changing question types insufficient without supervision
        \end{itemize}
        
        \item \textbf{Infrastructure is Primary Constraint}
        \begin{itemize}
            \item Technical failures most frequent complaint
            \item Client-side AI processing crucial for unstable connections
        \end{itemize}
        
        \item \textbf{User Experience Impacts Performance}
        \begin{itemize}
            \item Complex interfaces increase anxiety and errors
            \item Intuitive portals with minimal setup required
        \end{itemize}
    \end{enumerate}
\end{frame}

\begin{frame}{Lessons Learned (continued)}
    \begin{enumerate}
        \setcounter{enumi}{5}
        \item \textbf{LLM Consistency}
        \begin{itemize}
            \item Reliable, generalizable scoring across large datasets
            \item Outperforms traditional methods in efficiency
        \end{itemize}
        
        \item \textbf{Human-AI Collaboration}
        \begin{itemize}
            \item Dual-process models enhance accuracy
            \item LLMs as assistants, not standalone tools
        \end{itemize}
        
        \item \textbf{Prompting Techniques}
        \begin{itemize}
            \item Zero-shot and chain-of-thought prompting improve rubric-based evaluation
            \item Better handling of subjective criteria
        \end{itemize}
    \end{enumerate}
    
    \vspace{0.5cm}
    \begin{center}
        \textbf{These lessons directly inform Veritas design decisions}
    \end{center}
\end{frame}
